% Auto-generated from documents/thesis/05_conclusion.md by src.analysis.thesis_latex.
% Needs \usepackage{graphicx,booktabs} and your bibliography for \cite.

\chapter{Conclusion}

\section{Summary}

This thesis studied four-class ordinal student-engagement recognition from OpenFace facial descriptors and I3D motion features, and contributed a \textbf{semantic-group hybrid temporal architecture} together with a \textbf{cross-dataset generalization study}.

The central methodological idea --- decomposing the 709-dimensional OpenFace descriptor into five behavioral groups (gaze, eye, face, head pose, action units) and encoding each with its own temporal encoder, optionally fused with an I3D stream --- was validated by a systematic ablation over all 32 per-group encoder assignments. Three findings stand out:

1. \textbf{The I3D-fused hybrid is the best in-domain model and is robustly better than the strongest monolithic baseline.} The whole I3D-fused family sits above the baseline QWK of 0.537, and the best configuration (all-TCN) reaches \textbf{QWK 0.574} (accuracy 0.776). The gain over the best baseline is genuine but \textbf{incremental} (+0.037 QWK); the value lies in the framework and the systematic evidence, not a SOTA leap. 2. \textbf{Most groups are encoder-agnostic; only head pose clearly prefers a TCN} (+0.025 QWK), an interpretable design rule reflecting the temporally-local nature of head dynamics. 3. \textbf{Engagement classifiers do not generalize across datasets.} Off-diagonal QWK collapses to near zero (CMOSE$\rightarrow$DAiSEE 0.02) even when accuracy looks acceptable, so ordinal agreement metrics --- not accuracy --- must be reported on imbalanced engagement data. 4. \textbf{On a self-collected, hand-labeled private set --- the most realistic test --- the two contributions compound.} The private set is test-only, so the only lever is the training source; training on the \textbf{combined} corpus with the \textbf{semantic-group hybrid} gives the best result (QWK 0.365), beating the best baseline by +0.080 (vs +0.037 in-domain) and far exceeding any single-source training (CMOSE-only 0.19, DAiSEE-only 0.11). This is the practical payoff of the thesis. The rare Highly-Disengage class remains the principal open problem under distribution shift.

A note on scope: the architecture ablation is anchored on CMOSE because DAiSEE's in-domain ordinal signal is too weak (QWK 0.139) for architecture differences to be distinguishable from label noise; CMOSE is used for development and the private set for the decisive validation.

Practically, the thesis argues for two changes in how engagement models are built and reported: prefer \textbf{balanced/ordinal losses and QWK/macro-accuracy} over accuracy, and \textbf{pool source corpora} when the deployment distribution is unknown.

\section{Suggestion for Future Works}

\begin{itemize}
  \item \textbf{Targeting the rarest class.} Combine the architecture with imbalance-specific remedies (focal/ordinal losses tuned for HD, minority oversampling such as SMOTE \cite{smote}, or cost-sensitive thresholds) to recover Highly-Disengage recall without trading away HE.
  \item \textbf{Domain adaptation.} The cross-dataset gap invites explicit unsupervised domain adaptation or domain-invariant representation learning, rather than naive corpus pooling.
  \item \textbf{Learned, not enumerated, group encoders.} Replace the discrete 32-config search with a differentiable architecture search or a gating mechanism that selects per-group temporal operators end-to-end.
  \item \textbf{Frequency-domain analysis of engagement (exploratory).} A natural complementary direction is to analyze the \emph{temporal frequency} of facial signals --- the hypothesis that engaged behavior is low-frequency/stable while disengaged behavior is high-frequency/fidgety --- via spectral (e.g. STFT-based) features or convolutions over the frequency axis. This was not implemented here and is left as future work.
  \item \textbf{Multimodal extension.} CMOSE also provides audio/speech features; fusing them as additional semantic streams is a direct extension of the hybrid.
\end{itemize}

% --- note from draft (Markdown blockquote) ---
% TODO: tighten to your program's expected length and add any committee-requested directions.
